% ============================================================================
% Fichier  : partie1/P01_C04_methodologies.tex
% Titre    : Méthodologies et Référentiels
% Auteur   : MINKA MI NGUIDJOÏ Thierry Emmanuel
% Version  : 1.0 -- novembre 2026
% ============================================================================
% Sources : NIST SP 800-115 (2008) ; PTES ; OSSTMM 3 (ISECOM) ; OWASP WSTG ;
% MITRE ATT&CK ; FIRST, CVSS.
% ============================================================================

\chapter{Méthodologies et Référentiels}
\label{chap:methodes}
\chaptermark{Méthodologies et référentiels}

\epigraph{%
    « Une méthode ne remplace pas le talent~;\\
    elle empêche d'oublier ce que le talent n'a pas vu. »%
}{Principe du cours}

\niveauChapitre{Intermédiaire}{Chapitres~\ref{chap:ethique}
à~\ref{chap:engagement}.}

\begin{boxobjectifs}
\begin{itemize}
    \item Décrire les phases d'un test selon le NIST SP~800-115 et le PTES.
    \item Situer l'OSSTMM, l'OWASP WSTG et les autres guides.
    \item Utiliser MITRE ATT\&CK pour décrire un comportement d'attaquant.
    \item Exprimer la gravité d'une faille par le CVSS.
    \item Tenir le journal de mission qui rend un test reproductible.
\end{itemize}
\end{boxobjectifs}

\bigskip

% ============================================================================
\section*{Question directrice : comment être sûr de ne rien avoir oublié~?}
% ============================================================================

Au premier jour, M.~BIKORO remet à son équipe une liste de vérification
fondée sur des référentiels publics, et un journal de mission vierge.
Chaque test y sera noté. Quand LiZbiZ demandera « avez-vous testé la
messagerie~? », la réponse sera écrite, avec l'heure.

\medskip
\textbf{Plan du chapitre.}
\begin{enumerate}
    \item Les phases d'un test (section~\ref{sec:c04-phases}).
    \item Les guides spécialisés (section~\ref{sec:c04-guides}).
    \item Décrire l'attaquant~: ATT\&CK (section~\ref{sec:c04-attack}).
    \item Coter la gravité~: CVSS (section~\ref{sec:c04-cvss}).
    \item Le journal de mission (section~\ref{sec:c04-journal}).
\end{enumerate}

\separateur

% ============================================================================
\section{Les Phases d'un Test}
\label{sec:c04-phases}
% ============================================================================

\begin{tableaucours}{L{3.0cm}L{4.6cm}L{4.6cm}}{%
    \tableauHeader{Phase} &
    \tableauHeader{NIST SP~800-115} &
    \tableauHeader{PTES}
}
Préparer & Planification & Pré-engagement \\
Reconnaître & Découverte & Collecte de renseignements, modélisation des
menaces \\
Analyser & Découverte (vulnérabilités) & Analyse des vulnérabilités \\
Exploiter & Attaque & Exploitation, post-exploitation \\
Restituer & Rapport & Rapport \\
\end{tableaucours}
\vspace{-0.8em}
\begin{center}{\small\itshape Tableau~4.1. Les phases d'un test (d'après le NIST SP~800-115 et le PTES)}\end{center}

Le NIST SP~800-115 décrit les techniques de test et d'évaluation de la
sécurité de l'information~\cite{nist800115}. Le PTES (\emph{Penetration
Testing Execution Standard}) détaille sept phases, du pré-engagement au
rapport~\cite{ptes}. Les phases ne sont pas strictement successives~: une
découverte en exploitation relance la reconnaissance.

% ============================================================================
\section{Les Guides Spécialisés}
\label{sec:c04-guides}
% ============================================================================

\begin{tableaucours}{L{3.4cm}L{8.8cm}}{%
    \tableauHeader{Référentiel} &
    \tableauHeader{Apport}
}
OSSTMM (ISECOM) & Méthode de mesure de la sécurité opérationnelle, avec
des métriques de surface d'attaque~\cite{osstmm3} \\
OWASP WSTG & Guide de test des applications web, scénario par
scénario~\cite{owaspwstg} \\
OWASP MASTG et MASVS & Test et exigences des applications mobiles \\
NIST SP~800-115 & Cadre général, techniques, rapport \\
\end{tableaucours}
\vspace{-0.8em}
\begin{center}{\small\itshape Tableau~4.2. Guides de test}\end{center}

% ============================================================================
\section{Décrire l'Attaquant : ATT\&CK}
\label{sec:c04-attack}
% ============================================================================

\begin{boxdef}
\termecle{MITRE ATT\&CK}~: base de connaissances publique, fondée sur des
observations réelles, qui décrit les tactiques (le but de l'attaquant à une
étape) et les techniques (la manière de l'atteindre) des
adversaires~\cite{mitreattack}.
\end{boxdef}

\begin{tableaucours}{L{3.6cm}L{8.6cm}}{%
    \tableauHeader{Tactique} &
    \tableauHeader{Exemple dans l'Opération GAROUA}
}
Accès initial & Comptes valides volés par hameçonnage \\
Exécution & Lancement d'un outil sur le serveur de fichiers \\
Persistance & Création du compte svc\_maj2 \\
Élévation de privilèges & Exploitation d'un serveur non corrigé \\
Découverte & Balayage du réseau interne \\
Mouvement latéral & Connexions avec le mot de passe local partagé \\
Exfiltration & Copie de 40~Go vers un stockage externe \\
Impact & Chiffrement des serveurs \\
\end{tableaucours}
\vspace{-0.8em}
\begin{center}{\small\itshape Tableau~4.3. Quelques tactiques d'ATT\&CK appliquées à l'attaque de mars}\end{center}

Décrire un test en techniques ATT\&CK permet au SOC de vérifier, technique
par technique, s'il aurait détecté chaque geste (Chapitre~\ref{chap:rouge}).

% ============================================================================
\section{Coter la Gravité : CVSS}
\label{sec:c04-cvss}
% ============================================================================

\begin{boxdef}
\termecle{CVSS} (\emph{Common Vulnerability Scoring System}), publié par le
FIRST~: système qui exprime la gravité technique d'une vulnérabilité par un
score de 0 à 10, calculé à partir de métriques (vecteur d'attaque,
complexité, privilèges requis, interaction de l'utilisateur, impacts sur la
confidentialité, l'intégrité et la disponibilité)~\cite{first2019cvss31}.
\end{boxdef}

\begin{tableaucours}{L{3.0cm}L{2.6cm}}{%
    \tableauHeader{Score} &
    \tableauHeader{Qualification}
}
0,0 & Aucune \\
0,1 à 3,9 & Faible \\
4,0 à 6,9 & Moyenne \\
7,0 à 8,9 & Élevée \\
9,0 à 10,0 & Critique \\
\end{tableaucours}
\vspace{-0.8em}
\begin{center}{\small\itshape Tableau~4.4. Les qualifications CVSS}\end{center}

\begin{boxalert}
\textbf{Violation typique.} Classer les constats par le seul score CVSS.

\smallskip
\textbf{Règle~:} le score dit la gravité technique~; la priorité dépend
aussi de l'exposition, de la valeur de l'actif et de l'exploitation connue
(Chapitre~\ref{chap:vulnerabilites}). Le score de base se calcule, le
contexte se juge.
\end{boxalert}

% ============================================================================
\section{Le Journal de Mission}
\label{sec:c04-journal}
% ============================================================================

\begin{boxprocedure}
\textbf{Tenir le journal de mission}
\begin{enumerate}
    \item Une ligne par action~: heure, testeur, adresse source, cible,
          outil ou commande, résultat.
    \item Les captures et les fichiers de preuve, nommés et horodatés.
    \item Les incidents et les arrêts, avec qui a été prévenu.
    \item Les tests prévus mais non faits, et pourquoi.
    \item Conservation selon le contrat, accès restreint.
\end{enumerate}
\end{boxprocedure}

\begin{boxscenario}{Opération LOGONE. La question de la messagerie}
\textbf{Constatation.} Le 5~novembre, M.~KOUNOU demande si la messagerie en
ligne a été testée. Le journal répond~: test de la configuration
(authentification, règles de transfert) le 22~octobre de 9~h~10 à 11~h~40,
aucune faille~; pas de test de l'interface d'administration, hors périmètre.

\medskip
\textbf{Analyse} (confiance élevée). Sans journal, la réponse aurait été un
souvenir. Avec lui, c'est une constatation vérifiable.
\end{boxscenario}

\begin{boxmetier}{Pentesteur}
\textbf{Écrire au fil de l'eau.} Le journal se tient pendant le test, pas
le soir. Une action non notée est une action que l'on ne pourra ni
expliquer au client, ni défendre en cas d'incident.
\end{boxmetier}

\separateur

\begin{boxresume}
\textbf{Ce qu'il faut retenir de ce chapitre~:}
\begin{itemize}
    \item \textbf{Préparer, reconnaître, analyser, exploiter, restituer}~:
          les phases du NIST SP~800-115 et du PTES.
    \item \textbf{OSSTMM, OWASP WSTG, MASTG} spécialisent la méthode.
    \item \textbf{ATT\&CK} décrit l'attaquant en tactiques et techniques.
    \item \textbf{Le CVSS} cote la gravité technique~; la priorité y ajoute
          le contexte.
    \item \textbf{Le journal de mission} rend le test vérifiable.
\end{itemize}
\end{boxresume}

\begin{boxinfo}
\textbf{Points de contrôle de l'élève.} Je sais~: nommer les phases d'un
test~; choisir un guide spécialisé~; décrire une attaque en tactiques
ATT\&CK~; qualifier un score CVSS~; tenir un journal de mission.
\end{boxinfo}

\bigskip

\begin{boxexo}[title={\ding{45}\quad Exercice~4.1. Les phases \niveaubase}]
Dans quelle phase se range~: la lettre d'autorisation~; le balayage des
ports~; la démonstration d'accès aux administrateurs du domaine~?
\end{boxexo}

\begin{boxexo}[title={\ding{45}\quad Exercice~4.2. Qualifier \niveaubase}]
Qualifiez les scores CVSS 9,8~; 7,5~; 5,3~; 3,1.
\end{boxexo}

\begin{boxexo}[title={\ding{45}\quad Exercice~4.3. Tactiques \niveaubase}]
Rangez dans les tactiques d'ATT\&CK~: la suppression des sauvegardes~; la
connexion avec le mot de passe du technicien~; la copie de 40~Go.
\end{boxexo}

\begin{boxexo}[title={\ding{45}\quad Exercice~4.4. Le guide \niveaumoyen}]
Quel guide utiliser pour tester l'API de l'ANPP~? Et son application mobile~?
\end{boxexo}

\begin{boxexo}[title={\ding{45}\quad Exercice~4.5. Le contexte \niveaumoyen}]
Deux failles ont un score de 8,1~: l'une sur un serveur exposé à Internet,
l'autre sur un poste isolé. Laquelle corriger d'abord~? Pourquoi~?
\end{boxexo}

\begin{boxexo}[title={\ding{45}\quad Exercice~4.6. Le journal \niveaumoyen}]
Rédigez cinq lignes du journal de mission du 26~octobre (tests internes).
\end{boxexo}

\begin{boxexo}[title={\ding{45}\quad Exercice~4.7. Non fait \niveaumoyen}]
Pourquoi noter les tests prévus mais non faits~? Donnez deux raisons
légitimes de ne pas faire un test prévu.
\end{boxexo}

\begin{boxexo}[title={\ding{45}\quad Exercice~4.8. ATT\&CK et SOC \niveaumoyen}]
Comment la description d'un test en techniques ATT\&CK aide-t-elle le SOC
à s'améliorer~?
\end{boxexo}

\begin{boxexo}[title={\ding{45}\quad Exercice~4.9. Liste de vérification \niveauexpert}]
Construisez la liste de vérification des tests internes de l'Opération
LOGONE, phase par phase.
\end{boxexo}

\begin{boxexo}[title={\ding{45}\quad Exercice~4.10. Comparer \niveauexpert}]
Comparez le NIST SP~800-115, le PTES et l'OSSTMM~: objet, niveau de détail,
usage recommandé.
\end{boxexo}

\bigskip

\begin{boxinfo}
\textbf{Pour aller plus loin.}
\begin{itemize}
    \item NIST SP~800-115~\cite{nist800115}~; PTES~\cite{ptes}~; OSSTMM~\cite{osstmm3}.
    \item MITRE ATT\&CK~\cite{mitreattack}~; CVSS~v3.1~\cite{first2019cvss31}.
\end{itemize}
\end{boxinfo}

\begin{boxinfo}
\textbf{Transition.} La Partie~II commence par ce que l'on apprend sans
toucher aux systèmes~: la reconnaissance passive (Chapitre~\ref{chap:osint}).
\end{boxinfo}
